
\documentclass[12pt,a4paper]{article}

\usepackage[utf8]{inputenc}
\usepackage[T1]{fontenc}
\usepackage{graphicx}
\usepackage{amsmath}
\usepackage{amssymb}
\usepackage{geometry}
\usepackage{fancyhdr}
\usepackage{enumitem}
\usepackage{parskip}
\usepackage{xcolor}
\usepackage{anyfontsize}
\usepackage{subcaption}
\usepackage{booktabs}
\usepackage{array}
\usepackage{setspace}
\onehalfspacing
\usepackage[colorlinks=true]{hyperref}
\usepackage{xurl}
\usepackage{float}
\usepackage{longtable}
\usepackage{titlesec}
\usepackage{tabularx}
\usepackage{ragged2e}
\usepackage[most]{tcolorbox}
\usepackage{makecell}
\usepackage{listings}
\usepackage{caption}
\usepackage{pdflscape}
\usepackage{placeins}

\graphicspath{{Images/}}

% =====================================
% Colors and Styles
% =====================================
\definecolor{itiBlue}{RGB}{0,82,147}
\definecolor{absblack}{RGB}{0,0,0}
\definecolor{lightgray}{RGB}{245,245,245}
\definecolor{softblue}{RGB}{239,246,252}
\definecolor{softgreen}{RGB}{239,252,245}
\definecolor{darkgreen}{RGB}{33,110,57}
\definecolor{softorange}{RGB}{255,248,235}
\definecolor{darkorange}{RGB}{173,94,0}
\hypersetup{
    linkcolor = absblack,
    urlcolor = itiBlue,
    citecolor = itiBlue,
    colorlinks = true
}

% =====================================
% Page Layout
% =====================================
\geometry{a4paper, margin=0.75in}
\setlength{\parindent}{0pt}
\setlength{\parskip}{4pt}
\setlist[itemize]{leftmargin=1.0cm,itemsep=1pt,topsep=1pt}
\setlist[enumerate]{leftmargin=1.0cm,itemsep=1pt,topsep=1pt}
\setlength{\headheight}{15pt}
\setlength{\tabcolsep}{4pt}
\renewcommand{\arraystretch}{1.18}
\emergencystretch=3em
\sloppy

\pagestyle{fancy}
\fancyhf{}
\rhead{LTE Throughput Anomaly Detection Progress}
\lhead{DIGIS GP Team}
\renewcommand{\headrulewidth}{0.4pt}
\fancyfoot[C]{%
    \vspace{5pt}%
    \rule{\textwidth}{0.4pt}\\%
    \vspace{2pt}%
    \thepage%
}

\titleformat{\section}{\color{itiBlue}\normalfont\Large\bfseries}{\thesection}{0.75em}{}
\titleformat{\subsection}{\color{absblack}\normalfont\large\bfseries}{\thesubsection}{0.75em}{}
\titleformat{\subsubsection}{\color{absblack}\normalfont\normalsize\bfseries}{\thesubsubsection}{0.75em}{}

\captionsetup[table]{position=bottom,skip=5pt,font=small,labelfont=bf}
\captionsetup[figure]{position=bottom,skip=5pt,font=small,labelfont=bf}

\newcolumntype{Y}{>{\RaggedRight\arraybackslash}X}
\newcolumntype{L}[1]{>{\RaggedRight\arraybackslash}m{#1}}
\newcolumntype{P}[1]{>{\RaggedRight\arraybackslash}p{#1}}
\newcolumntype{C}[1]{>{\Centering\arraybackslash}p{#1}}

\tcbset{
    colback=softblue,
    colframe=itiBlue,
    boxrule=0.6pt,
    arc=2pt,
    left=6pt,
    right=6pt,
    top=6pt,
    bottom=6pt
}

\lstdefinestyle{cmdstyle}{
    basicstyle=\ttfamily\scriptsize,
    breaklines=true,
    columns=fullflexible,
    frame=single,
    rulecolor=\color{lightgray},
    backgroundcolor=\color{lightgray},
    xleftmargin=2pt,
    xrightmargin=2pt,
    aboveskip=6pt,
    belowskip=6pt,
    showstringspaces=false
}
\lstset{style=cmdstyle}

\newtcolorbox{infobox}[1]{
    title=#1,
    colback=softblue,
    colframe=itiBlue,
    fonttitle=\bfseries,
    boxrule=0.6pt,
    arc=2pt
}

\newtcolorbox{successbox}[1]{
    title=#1,
    colback=softgreen,
    colframe=darkgreen,
    fonttitle=\bfseries,
    boxrule=0.6pt,
    arc=2pt
}

\newtcolorbox{warningbox}[1]{
    title=#1,
    colback=softorange,
    colframe=darkorange,
    fonttitle=\bfseries,
    boxrule=0.6pt,
    arc=2pt
}

\newcommand{\fname}[1]{\texttt{#1}}
\newcommand{\colname}[1]{\texttt{#1}}
\newcommand{\flowarrow}{\ensuremath{\Longrightarrow}}

\begin{document}

\begin{titlepage}
    \centering
	\begin{figure}[H]
		\centering
		\IfFileExists{../ITI_Logo.png}{%
			
\includegraphics[width=0.68\linewidth]{../ITI_Logo}%
		}{%
			\fbox{\parbox[c][3cm][c]{0.7\linewidth}{\centering\Large ITI Logo\\[0.3cm]\normalsize Place \texttt{ITI\_Logo.png} in the required logo path}}
		}
	\end{figure}

    \vspace{0.9cm}
    \textcolor{absblack}{\rule{1\textwidth}{1.5pt}} \\[0.7cm]
    {\fontsize{21}{28}\selectfont\bfseries LTE Throughput Anomaly Detection}\\[0.35cm]
    {\fontsize{16}{18}\selectfont\bfseries Progress Report and Final Workflow Summary}\\[0.25cm]
    \textcolor{absblack}{\rule{1\textwidth}{1.5pt}} \\[1.1cm]

    \vspace*{0.2cm}
    \begin{tabular}{@{}l@{\hspace{1.3cm}}l@{}}
        \textbf{\fontsize{14}{16}\selectfont Created by:} & \fontsize{16}{18}\selectfont DIGIS GP Team \\ [0.2cm]
        & \fontsize{12}{14}\selectfont \textcolor{itiBlue}{\href{mailto:osamaeglal143@gmail.com}{osamaeglal143@gmail.com}} \\[0.2cm]
        & \fontsize{12}{14}\selectfont \textcolor{absblack}{Intake 46} \\[0.8cm]
        \textbf{\fontsize{14}{16}\selectfont Submitted to:} & \fontsize{16}{18}\selectfont Eng. Ibrahim El Shal \\[0.8cm]
        \textbf{\fontsize{14}{16}\selectfont Report Date:} & \fontsize{16}{18}\selectfont \today \\
    \end{tabular}
    \vfill
    \textcolor{absblack}{\rule{1\textwidth}{1pt}}
    \thispagestyle{empty}
\end{titlepage}

\small
\setstretch{1.05}
\tableofcontents
\newpage

\section{Executive Summary}

This report summarizes the progress made in the LTE downlink throughput anomaly-detection module. The work started from raw drive-test data and reached a validated, RCA-ready anomaly-detection workflow. The module does not perform final root-cause analysis (RCA). Instead, it detects suspicious LTE-DL throughput behavior, assigns interpretable anomaly types and scores, and exports structured evidence for the next RCA stage.

\begin{successbox}{Current outcome}
The latest workflow uses all valid LTE-DL samples, treats download activity as context only, uses RB allocation as the main demand/resource-confidence signal, and combines statistical anomaly methods with dual expected-throughput references: radio-condition P75 and RB-conditioned P75.
\end{successbox}

Main achievements:
\begin{itemize}
    \item Built a standalone notebook that starts from \fname{Network\_Drive\_Test.pkl} and creates cleaned, canonical LTE-DL modeling outputs.
    \item Validated telecom thresholds using throughput distributions and KPI boxplots.
    \item Improved the anomaly logic from simple low-throughput thresholds to context-aware detection using radio, RB, CA, event, and route evidence.
    \item Added episode-level grouping to reduce repeated adjacent anomaly rows and make the output more RCA-friendly.
    \item Prepared the path for a future nonlinear ML expected-throughput model.
\end{itemize}

\section{Project Scope and Module Boundary}

The full project is split into three layers:
\begin{enumerate}
    \item \textbf{Anomaly detection:} detect LTE throughput degradation and prepare evidence.
    \item \textbf{RCA:} use the anomaly evidence to decide likely technical causes.
    \item \textbf{LLM output layer:} connect anomaly and cause into a structured final output.
\end{enumerate}

This report covers the first layer only. Final RCA decisions are intentionally not assigned in the anomaly notebook. Labels such as \fname{radio\_limited\_degradation} or \fname{unexpected\_underperformance} are anomaly cases and evidence directions, not final root causes.

\begin{infobox}{Final anomaly-detection objective}
Detect low or degraded LTE-DL throughput, separate simple low-throughput cases from unexpected underperformance, and export a clean handoff table that includes throughput, radio, RB, CA, event, session, route, and score evidence.
\end{infobox}

\section{End-to-End Development Flow}

The notebook evolved through multiple review and calibration cycles. The final flow is shown below.

\begin{center}
\begin{tabularx}{\textwidth}{L{0.25\textwidth}Y}
\toprule
\textbf{Stage} & \textbf{Progress made} \\
\midrule
Raw loading and cleaning & Load the raw pickle, protect important LTE/location/event/activity columns, remove unusable sparse/constant columns, and export cleaning reports. \\
Canonical mapping & Map long vendor column names into stable aliases such as \colname{actual\_lte\_dl\_throughput}, \colname{lte\_sinr}, \colname{lte\_rsrp}, \colname{serving\_pci}, and \colname{rb\_usage\_value}. \\
LTE-DL selection & Keep valid LTE downlink throughput samples with LTE measurements. The latest version uses all valid LTE-DL rows, not only active-download windows. \\
Telecom knob validation & Validate low-throughput threshold, radio bins, RB demand buckets, CA behavior, and scoring thresholds using plots. \\
Statistical AD methods & Add fixed threshold, global/session percentiles, robust z-score, IQR, rolling drop, sudden drop, sustained degradation, P75 underperformance, and RB-efficiency checks. \\
Score and severity & Convert correlated flags into independent trigger families, combine base score, magnitude score, evidence bonus, and uncertainty caps. \\
RCA handoff & Export row-level anomaly table and episode-level summary with KPI evidence and context fields. \\
\bottomrule
\end{tabularx}
\end{center}

\section{Pre-Anomaly Throughput Understanding}

The first important validation was the LTE-DL throughput distribution before anomaly detection. This graph was used to choose the fixed low-throughput thresholds.

\begin{figure}[H]
    \centering
    \includegraphics[width=0.92\linewidth]{01_pre_anomaly_throughput_distribution.png}
    \caption{LTE-DL throughput distribution before anomaly detection. The graph supports using 5 Mbps as a very-low threshold and 10 Mbps as the main low-throughput threshold.}
\end{figure}

Key observations:
\begin{itemize}
    \item The global P10 is around 5 Mbps, so 5 Mbps represents the lower tail of the dataset.
    \item The 10 Mbps threshold is above the extreme lower tail but still clearly below the median throughput.
    \item Global P25 is around 15.43 Mbps and the median is around 38.80 Mbps, so using 15--20 Mbps as a hard low-throughput threshold would over-broaden the anomaly definition.
\end{itemize}

\begin{center}
\begin{tabularx}{\textwidth}{L{0.32\textwidth}L{0.18\textwidth}Y}
\toprule
\textbf{Knob} & \textbf{Value} & \textbf{Reason} \\
\midrule
Very-low throughput & 5 Mbps & Close to global P10, used for extreme low-throughput severity. \\
Low-throughput threshold & 10 Mbps & Practical low-throughput flag; strict enough not to cover too much of the distribution. \\
Soft low-throughput severity start & 20 Mbps & Used in scoring to gradually penalize samples below 20 Mbps, while keeping 10 Mbps as the hard flag. \\
\bottomrule
\end{tabularx}
\end{center}

\section{Radio-KPI Knob Validation}

Radio conditions were validated using throughput boxplots by RSRP, SINR, and RSRQ bins. These plots were used to decide whether the bins follow meaningful throughput behavior.

\subsection{RSRP and SINR}

\begin{figure}[H]
    \centering
    \begin{subfigure}{0.49\linewidth}
        \centering
        \includegraphics[width=\linewidth]{02_throughput_by_rsrp_bin.png}
        \caption{Throughput by RSRP bin.}
    \end{subfigure}
    \hfill
    \begin{subfigure}{0.49\linewidth}
        \centering
        \includegraphics[width=\linewidth]{04_throughput_by_sinr_bin.png}
        \caption{Throughput by SINR bin.}
    \end{subfigure}
    \caption{RSRP and SINR show clear monotonic improvement in throughput, supporting their use as strong radio-quality knobs.}
\end{figure}

The final RSRP bins are:
\begin{center}
\begin{tabular}{ll}
\toprule
\textbf{RSRP range} & \textbf{Bin} \\
\midrule
$<$ -120 dBm & Poor \\
$[-120, -100)$ dBm & Fair \\
$[-100, -90)$ dBm & Good \\
$\geq$ -90 dBm & Excellent \\
\bottomrule
\end{tabular}
\end{center}

The final SINR bins are:
\begin{center}
\begin{tabular}{ll}
\toprule
\textbf{SINR range} & \textbf{Bin} \\
\midrule
$<$ 0 dB & Poor \\
$[0, 10)$ dB & Fair \\
$[10, 15)$ dB & Good \\
$\geq$ 15 dB & Excellent \\
\bottomrule
\end{tabular}
\end{center}

\subsection{RSRQ}

\begin{figure}[H]
    \centering
    \includegraphics[width=0.72\linewidth]{03_throughput_by_rsrq_bin.png}
    \caption{Throughput by RSRQ bin. RSRQ is useful as evidence, but it is less monotonic than RSRP and SINR.}
\end{figure}

The final RSRQ bins are:
\begin{center}
\begin{tabular}{ll}
\toprule
\textbf{RSRQ range} & \textbf{Bin} \\
\midrule
$<$ -20 dB & Poor \\
$[-20, -15)$ dB & Fair \\
$[-15, -10)$ dB & Good \\
$\geq$ -10 dB & Excellent \\
\bottomrule
\end{tabular}
\end{center}

RSRQ was kept as supporting radio evidence, but not as the strongest standalone trigger because its plotted relationship with throughput is less stable.

\section{Demand and Resource Validation Using RB}

A major improvement was replacing activity-window filtering with RB-based demand/resource confidence. The throughput distributions showed that active-download and outside-download contexts were close enough that activity context should not define the main analysis population. RB usage is more directly related to what the scheduler allocated.

\begin{figure}[H]
    \centering
    \includegraphics[width=0.80\linewidth]{05e_throughput_by_rb_usage_bucket.png}
    \caption{Throughput by RB usage bucket. The monotonic increase from low RB to high RB supports using RB as the main demand/resource-confidence signal.}
\end{figure}

Final RB interpretation:
\begin{itemize}
    \item \textbf{Low RB + low throughput:} uncertain; may be low demand, low scheduling, or resource limitation.
    \item \textbf{Medium/high RB + low throughput:} much more suspicious because resources were allocated but throughput stayed poor.
    \item \textbf{High RB + good radio + low throughput:} strong unexpected-underperformance candidate.
\end{itemize}

Carrier aggregation was also validated as an important throughput knob.

\begin{figure}[H]
    \centering
    \includegraphics[width=0.72\linewidth]{05_throughput_by_carrier_count.png}
    \caption{Throughput by carrier count. Higher carrier count generally increases achievable throughput, supporting CA-aware anomaly logic.}
\end{figure}

\section{Expected Throughput Design}

A central part of the work was the expected-throughput design. Early P75 logic was useful but could be too broad or optimistic. The final design uses two complementary expected-throughput references.

\begin{center}
\begin{tabularx}{\textwidth}{L{0.28\textwidth}L{0.25\textwidth}Y}
\toprule
\textbf{Expected reference} & \textbf{Includes RB?} & \textbf{Purpose} \\
\midrule
Radio P75 & No & Estimates what throughput is normally achievable under similar radio and CA conditions. It can reveal scheduler/resource/demand underperformance. \\
RB-conditioned P75 & Yes & Estimates what throughput is normally achievable after also considering RB allocation. It is stronger evidence of throughput inefficiency. \\
\bottomrule
\end{tabularx}
\end{center}

Important assumptions:
\begin{itemize}
    \item Minimum P75 group size: 30 samples.
    \item Strict P75 anomaly decisions use high-confidence fallback levels only.
    \item Strong underperformance requires expected throughput to be meaningful, a large expected-actual gap, a low actual/expected ratio, and medium/high RB demand when RB is reliable.
\end{itemize}

The current strict underperformance logic is summarized as:
\begin{lstlisting}
Expected TP is high enough
AND actual / expected is low
AND expected - actual gap is large
AND P75 confidence is high
AND RB demand is medium/high
\end{lstlisting}

\section{Statistical Anomaly Detection Methods}

The notebook combines several methods because no single method captures all LTE degradation patterns.

\begin{center}
\begin{tabularx}{\textwidth}{L{0.28\textwidth}Y}
\toprule
\textbf{Method} & \textbf{Role in the final detector} \\
\midrule
Fixed low throughput & Flags samples below 10 Mbps. \\
Global P10 & Flags bottom-tail samples relative to the full LTE-DL distribution. \\
Session P10 & Detects samples that are poor relative to their own drive-test session. \\
Robust z-score & Detects robust low outliers on log throughput. \\
IQR lower fence & Adds another robust distribution-based low-tail check. \\
Rolling local drop & Detects local dips compared with nearby samples. \\
Sudden drop & Detects abrupt one-sample drops from previous throughput. \\
Sustained window & Detects continuous low-throughput periods, not just sudden changes. \\
Radio/RB P75 & Detects underperformance against expected achievable throughput. \\
RB efficiency & Detects low throughput per RB when RB demand is medium/high. \\
\bottomrule
\end{tabularx}
\end{center}

Important method knobs:
\begin{itemize}
    \item Rolling window: 11 samples, kept sample-based because sample gaps were sufficiently consecutive.
    \item Sudden-drop rule: previous throughput at least 10 Mbps, current/previous ratio at most 0.5, and gap at least 5 Mbps.
    \item Robust z-score threshold: approximately -3.0 for lower-tail outliers.
    \item CA underperformance: CA active, medium/high RB, expected radio P75 at least about 40 Mbps, actual throughput around 15 Mbps or less, and no obvious poor RSRP/SINR.
    \item Event windows: any event $\pm$30 s, handover $\pm$5 s, RACH failure $\pm$10 s, CA change $\pm$10 s.
\end{itemize}

\section{Score Combination and Final Calibration}

The scoring layer was improved after observing that simple flag stacking and fixed floors can distort the score distribution. The final scoring design uses independent trigger families and smooth magnitude terms.

\begin{figure}[H]
    \centering
    \includegraphics[width=0.78\linewidth]{11_anomaly_method_percentages.png}
    \caption{Anomaly method and case flagging rate. The chart is used to check whether one method dominates the final anomaly decision.}
\end{figure}

Trigger families prevent correlated methods from over-stacking. For example, fixed low throughput, global P10, session P10, robust z-score, and IQR all belong to a low-throughput family because they mostly say the same thing: the throughput value is low.

The final score combines:
\begin{enumerate}
    \item Base score from the strongest trigger.
    \item Continuous magnitude score using ratio severity, gap severity, low-throughput severity, drop severity, and RB-efficiency severity.
    \item Small capped bonus for multiple independent trigger families.
    \item Supporting evidence bonus from RB, CA, radio, and events.
    \item Low-RB uncertainty cap and score saturation control.
\end{enumerate}

\begin{figure}[H]
    \centering
    \includegraphics[width=0.78\linewidth]{10b_anomaly_score_distribution_flagged_only.png}
    \caption{Score distribution for flagged samples only. This plot is used to evaluate whether the score is too compressed, too spiky, or too saturated.}
\end{figure}

Severity thresholds used for presentation are:
\begin{center}
\begin{tabular}{ll}
\toprule
\textbf{Score range} & \textbf{Severity} \\
\midrule
35--60 & Medium \\
60--80 & High \\
$\geq$ 80 & Critical \\
\bottomrule
\end{tabular}
\end{center}

\section{Spatial and Contextual Validation}

Route plots were added to check whether anomalies are spatially reasonable and not random isolated artifacts.

\begin{figure}[H]
    \centering
    \begin{subfigure}{0.49\linewidth}
        \centering
        \includegraphics[width=\linewidth]{09_route_by_throughput.png}
        \caption{Route colored by throughput.}
    \end{subfigure}
    \hfill
    \begin{subfigure}{0.49\linewidth}
        \centering
        \includegraphics[width=\linewidth]{17c_route_with_anomaly_overlay.png}
        \caption{Route with anomaly overlay.}
    \end{subfigure}
    \caption{Spatial sanity checks connect throughput behavior and anomaly locations to the driven route.}
\end{figure}

Context-window plots were also added to inspect selected anomaly episodes around the selected timestamp.

\begin{figure}[H]
    \centering
    \begin{subfigure}{0.49\linewidth}
        \centering
        \includegraphics[width=\linewidth]{18_top_anomaly_episode_context_rank_4_throughput_context.png}
        \caption{Throughput context around a selected anomaly.}
    \end{subfigure}
    \hfill
    \begin{subfigure}{0.49\linewidth}
        \centering
        \includegraphics[width=\linewidth]{18_top_anomaly_episode_context_rank_4_radio_context.png}
        \caption{Radio context around the same anomaly.}
    \end{subfigure}
    \caption{Example anomaly context window. These plots help verify whether the anomaly is sudden, sustained, radio-limited, or underperforming against expected throughput.}
\end{figure}

\section{Final Output Organization}

The latest notebook separates outputs into three folders:
\begin{center}
\begin{tabularx}{\textwidth}{L{0.34\textwidth}Y}
\toprule
\textbf{Folder} & \textbf{Purpose} \\
\midrule
\fname{01\_preprocessing\_cleaning/} & Cleaning reports, canonical mapping, data dictionaries, and preserved activity/event tables. \\
\fname{02\_anomaly\_detection/} & Row-level anomaly output, episode summary, method summary, P75 reference tables, score QC, and PCI summaries. \\
\fname{03\_plots/} & Validation plots for throughput distribution, radio bins, RB demand, route maps, score distribution, and anomaly context windows. \\
\bottomrule
\end{tabularx}
\end{center}

For RCA integration, the recommended primary table is the episode-level anomaly summary, because it groups repeated adjacent anomaly samples into more meaningful anomaly incidents. The row-level anomaly file remains useful for detailed drill-down.

\section{Main Lessons Learned}

\begin{successbox}{Key technical lessons}
\begin{itemize}
    \item RSRP and SINR are strong radio knobs because their throughput relationship is visibly monotonic.
    \item RSRQ is useful, but should be supporting evidence rather than a dominant standalone trigger.
    \item RB usage is the best available demand/resource-confidence signal in this dataset.
    \item Activity-window labels should remain context only because their throughput distributions were close.
    \item Low throughput under low RB should not be scored like low throughput under high RB.
    \item P75 expected throughput is useful, but it must be protected with group-size, fallback, and confidence logic.
    \item Episode-level handoff is better than sample-only handoff for RCA.
\end{itemize}
\end{successbox}

\section{Future Work: Machine Learning Expected Throughput}

The next recommended step is to add a nonlinear ML expected-throughput model. The reason is that LTE throughput is not linear. It depends on interacting radio, RB, CA, PHY, mobility, and event features.

Recommended ML direction:
\begin{lstlisting}
Train nonlinear regression model
Input: radio + RB + CA + event + session features
Target: actual_lte_dl_throughput
Output: ml_expected_tp
Anomaly: actual_tp << ml_expected_tp
\end{lstlisting}

Proposed phases:
\begin{enumerate}
    \item \textbf{ML-1: Mean expected throughput.} Train a HistGradientBoostingRegressor or RandomForestRegressor to predict expected mean throughput.
    \item \textbf{ML-2: P75 expected throughput.} Train a quantile GradientBoostingRegressor with alpha = 0.75 to create a nonlinear version of the P75 expected-throughput method.
    \item \textbf{ML-3: Hybrid scoring.} Add ML expected-throughput underperformance as a new independent trigger family in the existing score.
    \item \textbf{ML-4: Explainability.} Add feature importance or SHAP explanations so the RCA team can understand why the model expected higher throughput.
\end{enumerate}

Important ML controls:
\begin{itemize}
    \item Use session-based train/test split, not random row split.
    \item Avoid leakage from throughput-derived columns such as PCell throughput, SCell throughput, activity throughput, expected P75, residuals, and existing anomaly flags.
    \item Train first on high-confidence normal/reference rows, then compare against the statistical detector.
    \item Keep statistical flags as explainable backup evidence even after adding ML.
\end{itemize}

\begin{infobox}{Final direction}
The most useful ML extension is not a black-box anomaly classifier. It is a nonlinear expected-throughput regression model that estimates what throughput should have been achievable, then flags large negative residuals under meaningful RB demand.
\end{infobox}

\end{document}
